% SPDX-License-Identifier: MIT
% Copyright (c) 2026 Evgeny Baulin
% Optimization Methods, Seminar 08: theory handout (Russian).
% Generated by checks/08. Constrained Optimality and Discrete Modelling.ipynb from web/08/js/content-*.js and content/students.json. Do not edit by hand.
% Compile with: latexmk -xelatex -interaction=nonstopmode -halt-on-error Theory_ru.tex
\documentclass[11pt,a4paper]{article}

\usepackage{fontspec}
\setmainfont{cmunrm.otf}[BoldFont=cmunbx.otf, ItalicFont=cmunti.otf, BoldItalicFont=cmunbi.otf]
\setsansfont{cmunss.otf}[BoldFont=cmunsx.otf, ItalicFont=cmunsi.otf, BoldItalicFont=cmunso.otf]
\setmonofont{cmuntt.otf}[BoldFont=cmuntb.otf, ItalicFont=cmunit.otf, BoldItalicFont=cmuntx.otf]
\usepackage{polyglossia}
\setdefaultlanguage{russian}
\setotherlanguage{english}
\newfontfamily\cyrillicfont{cmunrm.otf}[Ligatures=TeX, BoldFont=cmunbx.otf, ItalicFont=cmunti.otf, BoldItalicFont=cmunbi.otf]
\newfontfamily\cyrillicfontsf{cmunss.otf}[Ligatures=TeX, BoldFont=cmunsx.otf, ItalicFont=cmunsi.otf, BoldItalicFont=cmunso.otf]
\newfontfamily\cyrillicfonttt{cmuntt.otf}[BoldFont=cmuntb.otf, ItalicFont=cmunit.otf, BoldItalicFont=cmuntx.otf]

\usepackage[a4paper,left=2.3cm,right=2.3cm,top=2.4cm,bottom=2.4cm,headheight=14pt]{geometry}
\usepackage{amsmath,amssymb}
\usepackage{graphicx}
\usepackage{array,booktabs,longtable}
\usepackage{enumitem}
\usepackage{xcolor}
\usepackage[most]{tcolorbox}
\usepackage{titlesec}
\usepackage{fancyhdr}
\usepackage{needspace}
\usepackage[expansion=false]{microtype}
\usepackage[colorlinks=true,linkcolor=hseblue,urlcolor=hseblue,citecolor=hseblue]{hyperref}

\definecolor{hseblue}{RGB}{16,45,105}
\definecolor{hseaccent}{RGB}{198,40,40}
\definecolor{hsegray}{RGB}{108,114,126}
\definecolor{hselight}{RGB}{240,243,248}
\definecolor{hsewarm}{RGB}{253,244,232}

\hypersetup{pdftitle={Методы оптимизации. Семинар 8. Конспект теории}, pdfauthor={Евгений Баулин}}

\titleformat{\section}{\color{hseblue}\Large\bfseries}{\thesection}{0.6em}{}
\titleformat{\subsection}{\color{hseblue}\large\bfseries}{\thesubsection}{0.6em}{}
\titlespacing*{\section}{0pt}{16pt plus 4pt minus 2pt}{8pt plus 2pt}
\titlespacing*{\subsection}{0pt}{12pt plus 3pt minus 2pt}{6pt plus 2pt}
\pretocmd{\section}{\Needspace{7\baselineskip}}{}{\errmessage{cannot patch section}}
\pretocmd{\subsection}{\Needspace{5\baselineskip}}{}{\errmessage{cannot patch subsection}}

\tcbset{handoutbox/.style={enhanced, breakable, sharp corners, boxrule=0pt, left=7pt, right=7pt, top=5pt, bottom=5pt, before skip=8pt plus 2pt, after skip=8pt plus 2pt}}
\newtcolorbox{examplebox}{handoutbox, colback=hselight, colframe=hselight, before upper={\textcolor{hseblue}{\bfseries Пример.}\ }}
\newtcolorbox{whybox}{handoutbox, colback=white, colframe=hseblue, leftrule=2.5pt, top=3pt, bottom=3pt, before upper={\textcolor{hseaccent}{\bfseries Почему.}\ }}
\newtcolorbox{answerbox}{handoutbox, colback=hsewarm, colframe=hsewarm, before upper={\textcolor{hseaccent}{\bfseries Ответ.}\ }}
\newtcolorbox{conventionbox}{handoutbox, colback=hselight, colframe=hseblue, leftrule=2.5pt, before upper={\textcolor{hseblue}{\bfseries Обозначения.}\par\smallskip}}

\newcommand{\handoutfigure}[3]{\par\medskip\noindent\begin{minipage}{\linewidth}\centering
  \includegraphics[width=#1\linewidth]{figures/#2.pdf}\par\smallskip
  \begin{minipage}{0.94\linewidth}\small\textcolor{hseblue}{\bfseries Рисунок.}\ #3\end{minipage}\end{minipage}\par\medskip}
\newcommand{\problemhead}[3]{\par\bigskip\noindent{\color{hseblue}\large\bfseries #1.\ #2}\par\nopagebreak
  \noindent{\small\itshape\color{hsegray}#3}\par\nopagebreak\smallskip}
\newtcolorbox{definitionbox}[1]{handoutbox, colback=white, colframe=hseblue, leftrule=2.5pt, before upper={\textcolor{hseblue}{\bfseries #1}\ }}
\newtcolorbox{theorembox}[1]{handoutbox, colback=hselight, colframe=hseblue, leftrule=2.5pt, before upper={\textcolor{hseblue}{\bfseries #1}\ }}
\newtcolorbox{proofbox}{handoutbox, colback=white, colframe=hsegray, leftrule=1.2pt, top=3pt, bottom=3pt, before upper={\textcolor{hsegray}{\bfseries\itshape Доказательство.}\ }}
\newtcolorbox{notebox}[1]{handoutbox, colback=hselight, colframe=hselight, before upper={\textcolor{hseblue}{\bfseries #1}\par\smallskip}}
\newtcolorbox{trapbox}{handoutbox, colback=hsewarm, colframe=hseaccent, leftrule=2.5pt, before upper={\textcolor{hseaccent}{\bfseries Ловушка.}\ }}
\newtcolorbox{selfcheckbox}{handoutbox, colback=white, colframe=hseaccent, leftrule=2.5pt, top=3pt, bottom=3pt, before upper={\textcolor{hseaccent}{\bfseries Проверьте себя.}\ }}
\newcommand{\handoutdisplay}[1]{\sbox0{$\displaystyle #1$}\ifdim\wd0>\linewidth
  \par{\centering\binoppenalty=10000\relpenalty=0\relax$\displaystyle #1$\par}\else\[#1\]\fi}
\newcommand{\selfcheckanswer}{\par\smallskip\textcolor{hsegray}{\bfseries Ответ.}\ }

\setlist[itemize]{itemsep=3pt,topsep=4pt}
\setlist[enumerate]{itemsep=3pt,topsep=4pt}
\setlength{\parindent}{0pt}
\setlength{\parskip}{5pt plus 1pt}
\setlength{\emergencystretch}{2em}
\setcounter{tocdepth}{2}

\pagestyle{fancy}
\fancyhf{}
\fancyhead[L]{\small\color{hsegray}Методы оптимизации}
\fancyhead[R]{\small\color{hsegray}Семинар 8}
\fancyfoot[L]{\small\color{hsegray}Евгений Баулин}
\fancyfoot[C]{\small\thepage}
\fancyfoot[R]{\small\color{hsegray}Лицензия MIT}
\renewcommand{\headrulewidth}{0pt}
\renewcommand{\footrulewidth}{0pt}
\fancypagestyle{plain}{\pagestyle{fancy}}

\begin{document}

\begin{center}
  {\large\color{hsegray}Методы оптимизации\par}
  \vspace{4pt}
  {\large\bfseries\color{hseaccent}Семинар 8\par}
  \vspace{8pt}
  {\LARGE\bfseries\color{hseblue}Условная оптимальность и дискретное моделирование: ККТ, двойственность, целочисленные модели\par}
  \vspace{8pt}
  {\Large Конспект теории\par}
  \vspace{12pt}
  НИУ ВШЭ, факультет компьютерных наук\par
  2026-2027 учебный год\par
  \vspace{8pt}
  Лектор курса: Андрей Игнатов\par
  Евгений Баулин\par
  Telegram: \href{https://t.me/tarakan_tuc}{@tarakan\_tuc}\par
\end{center}

\noindent{\color{hseblue}\rule{\linewidth}{0.6pt}}

\tableofcontents

\clearpage


Этот конспект сопровождает семинар 8. Он ведёт от условий первого порядка в задаче с ограничениями к двойственной задаче и смыслу множителей, а затем переходит к моделям с целочисленными переменными, где та же идея оценки возвращается в виде релаксации. Все девять задач занятия разобраны полностью.

\begin{conventionbox}
Задача: минимизировать $f(x)$ по $x\in\mathbb R^n$ при $g_i(x)\le0$ и $h_j(x)=0$. Множители $\lambda_i\ge0$ — у неравенств, $\nu_j$ любого знака — у равенств. $L=f+\sum_i\lambda_ig_i+\sum_j\nu_jh_j$, $q=\inf_xL$, $p^*$ и $d^*$ — оптимальные значения прямой и двойственной задач. В целочисленной части задача максимизации сохраняет свой знак: там релаксация оценивает оптимум сверху.
\end{conventionbox}

% =====================================================================
\section{Условия оптимальности при ограничениях}\label{sec:kkt}

В точке безусловного минимума $\nabla f=0$. С ограничениями это перестаёт быть верным, как только решение оказывается на границе: функция хотела бы двигаться дальше, но её останавливает стенка. Верным остаётся равновесие между тягой $f$ и упором тех стенок, которых точка касается.

\begin{definitionbox}{Определение: Активное множество.}
Для допустимого $x$ индексы $i$, при которых $g_i(x)=0$, образуют активное множество $A(x)$. Ограничения с $g_i(x)<0$ выполнены строго и малые сдвиги не стесняют.
\end{definitionbox}

\begin{theorembox}{Теорема: Условия Каруша–Куна–Таккера.}
Пусть $x^*$ — точка локального минимума и в ней выполнено условие регулярности (например, градиенты активных ограничений и всех $h_j$ линейно независимы, или задача выпукла и выполнено условие Слейтера). Тогда найдутся $\lambda^*$ и $\nu^*$, при которых $\nabla_xL(x^*,\lambda^*,\nu^*)=0$; $x^*$ допустима; $\lambda^*\ge0$; $\lambda_i^*g_i(x^*)=0$ для всех $i$.
\end{theorembox}

\begin{whybox}
Стационарность говорит, что $-\nabla f(x^*)=\sum_i\lambda_i^*\nabla g_i(x^*)+\sum_j\nu_j^*\nabla h_j(x^*)$: направление быстрейшего убывания $f$ смотрит наружу из допустимого множества, в конус внешних нормалей тех стенок, которых касается точка. Любой допустимый сдвиг тогда ведёт вверх.
\end{whybox}

\begin{theorembox}{Утверждение: Достаточность в выпуклом случае.}
Если $f$ и все $g_i$ выпуклы, а все $h_j$ аффинны, то допустимая точка, удовлетворяющая условиям ККТ с какими-нибудь множителями, — точка глобального минимума.
\end{theorembox}

\begin{proofbox}
При фиксированных $\lambda^*\ge0$ функция $x\mapsto L(x,\lambda^*,\nu^*)$ выпукла и имеет нулевой градиент в $x^*$, поэтому $x^*$ минимизирует её на всём $\mathbb R^n$. Дополняющая нежёсткость и $h(x^*)=0$ дают $L(x^*,\lambda^*,\nu^*)=f(x^*)$. Для допустимого $x$ добавленные слагаемые не больше нуля, значит, $f(x^*)\le L(x,\lambda^*,\nu^*)\le f(x)$.
\end{proofbox}

\begin{itemize}
  \item Запишите $L$ и уравнения стационарности.
  \item Выберите предполагаемое активное множество $A$. Положите $g_i=0$ при $i\in A$ и $\lambda_i=0$ при $i\notin A$ и решите получившуюся квадратную систему.
  \item Отбросьте кандидата, если какое-то $g_i>0$ или какое-то $\lambda_i<0$.
  \item В выпуклой задаче первый уцелевший кандидат и есть ответ; иначе сравните значения всех уцелевших и убедитесь, что минимум вообще существует.
\end{itemize}

\handoutfigure{0.62}{fig_a1}{В задаче A1 минимум лежит на прямой $g_1$; $-\nabla f$ — положительное кратное нормали к этой прямой.}

\problemhead{A1}{Четыре варианта активного множества}{У доски}
\textbf{Условие.} Минимизируйте $f(x,y)=(x-3)^2+2(y-2)^2$ при $g_1=x+y-2\le0$ и $g_2=-y\le0$. Переберите все четыре активных множества, найдите точку ККТ и её множители и укажите первую координату $x^*$, $\lambda_1^*$ и $f(x^*)$.

\textbf{Решение.}
\begingroup\sloppy
\begin{enumerate}[label=\textbf{\arabic*.}, leftmargin=*]
  \item Ни одно ограничение не активно: безусловный минимум $(3;\ 2)$ даёт $g_1=3>0$ — точка недопустима, и хотя бы одно ограничение выполняется как равенство.
  \item Активно только $g_2$: $y=0$, $\lambda_1=0$, и первое уравнение даёт $x=3$; в точке $(3;\ 0)$ имеем $g_1=1>0$ — отбрасываем. Активны оба: из $x+y=2$ и $y=0$ получаем $(2;\ 0)$, а стационарность даёт $\lambda=(2;\ -6)$. Отрицательный множитель означает, что ограничение $y\ge0$ толкает точку не в ту сторону, — отбрасываем.
  \item Активно только $g_1$, $\lambda_2=0$: $x=3-\lambda_1/2$, $y=2-\lambda_1/4$, и из $x+y=2$ получаем $5-3\lambda_1/4=2$, то есть $\lambda_1=4$. Кандидат — $x^*=(1;\ 1)$ с $\lambda^*=(4;\ 0)$; $y^*>0$, значит, $g_2$ действительно неактивно.
  \item В точке $x^*$ градиент $\nabla f=(-4;\ -4)=-\lambda_1^*\nabla g_1$: антиградиент смотрит прямо наружу через прямую $x+y=2$. Задача выпуклая, поэтому точка ККТ — глобальный минимум, и $f(x^*)=6$.
\end{enumerate}
\endgroup
\begin{answerbox}
$x^*=(1;\ 1)$, $\lambda^*=(4;\ 0)$, $f(x^*)=6$.
\end{answerbox}

\begin{trapbox}
Активное не значит с положительным множителем. Импликация работает в одну сторону: из $\lambda_i>0$ следует $g_i=0$, но при $g_i=0$ возможно $\lambda_i=0$. Типичный случай — ограничение, которое лишь касается безусловного решения.
\end{trapbox}

\problemhead{A2}{Три ограничения}{Самостоятельно}
\textbf{Условие.} Минимизируйте $f(x,y)=3(x-5)^2+(y-2)^2$ при $g_1=x-3{,}5\le0$, $g_2=x+y-5{,}5\le0$ и $g_3=-y\le0$. Найдите точку ККТ, активное множество и все три множителя. Укажите $\lambda_1^*$, $\lambda_2^*$ и $f(x^*)$.

\textbf{Решение.}
\begingroup\sloppy
\begin{enumerate}[label=\textbf{\arabic*.}, leftmargin=*]
  \item Без ограничений $f$ достигала бы минимума в $(5,2)$, но там $g_1>0$: граница $x\le3{,}5$ нарушена, значит, $g_1$ обязано быть точным. Положим $x=3{,}5$ и пока оставим остальные множители нулевыми. Уравнение по $y$, $2(y-2)=0$, даёт $y=2$, уравнение по $x$, $6(x-5)+\lambda_1=0$, даёт $\lambda_1=9$.
  \item В точке $x^*=(3{,}5;\ 2)$ значения ограничений равны $g=(0;\ 0;\ -2)$. Ограничение на сумму $g_2$ тоже обращается в равенство, хотя мы его об этом не просили, а $g_3$ выполнено строго, поэтому $\lambda_3=0$ по дополняющей нежёсткости. Активное множество — $\{1,2\}$.
  \item Когда в игре оба ограничения $g_1$ и $g_2$, стационарность выглядит так: $6(x-5)+\lambda_1+\lambda_2=0$ и $2(y-2)+\lambda_2-\lambda_3=0$. При $y=2$, $\lambda_3=0$ второе уравнение вынуждает $\lambda_2=0$, а первое снова даёт $\lambda_1=9$. Градиенты $(1,0)$ и $(1,1)$ точных ограничений линейно независимы, так что других множителей нет: $\lambda^*=(9;\ 0;\ 0)$.
  \item Все четыре условия выполнены, $f$ выпукла, ограничения линейны — значит, это глобальный минимум, $f^*=6{,}75$. Ограничение $g_2$ касается точки $x^*$, но ничего не весит: уберите его — ничего не сдвинется; уберите $g_1$ — ответ изменится.
\end{enumerate}
\endgroup
\begin{answerbox}
$x^*=(3{,}5;\ 2)$, $\lambda_1=9$, $\lambda_2=0$, $\lambda_3=0$, $f^*=6{,}75$.
\end{answerbox}

\problemhead{A3}{Линейная цель на круге}{Самостоятельно}
\textbf{Условие.} Минимизируйте $x+2y$ при $x^2+y^2\le45$. Объясните, почему минимум лежит на границе, найдите его из условий ККТ и проверьте, что он глобальный. Укажите $y^*$, $\lambda^*$ и $f(x^*)$.

\textbf{Решение.}
\begingroup\sloppy
\begin{enumerate}[label=\textbf{\arabic*.}, leftmargin=*]
  \item При $c=(1;\ 2)$ и $g=x^2+y^2-45$ стационарность имеет вид $c+2\lambda(x,y)=0$, откуда $(x,y)=-c/(2\lambda)$: минимум лежит напротив $c$, и нужно $\lambda>0$.
  \item Окружность точна: $\|c\|^2/(4\lambda^2)=45$, где $\|c\|^2=5$, поэтому $\lambda^*=0{,}1667$, $x^*=(-3;\ -6)$ и $f^*=-15$.
  \item Цель линейна, круг выпуклый, значит, эта точка ККТ — глобальный минимум; вторая точка окружности на той же прямой — максимум, и множитель у неё отрицательный.
\end{enumerate}
\endgroup
\begin{answerbox}
$x^*=(-3;\ -6)$, $\lambda^*=0{,}1667$, $f^*=-15$.
\end{answerbox}

\begin{selfcheckbox}
Почему в задаче A1 случай с двумя активными ограничениями отбрасывается, хотя его точка допустима?
\selfcheckanswer Множитель при $y\ge0$ получается отрицательным, $\lambda_2=-6$: стенке $y=0$ пришлось бы тянуть точку к себе, а не отталкивать.
\end{selfcheckbox}

% =====================================================================
\section{Двойственность Лагранжа и цены}\label{sec:duality}

\begin{definitionbox}{Определение: Двойственная функция и двойственная задача.}
$q(\lambda,\nu)=\inf_xL(x,\lambda,\nu)$, возможно $-\infty$. Двойственная задача — максимизировать $q$ по $\lambda\ge0$ и всем $\nu$; её значение — $d^*$.
\end{definitionbox}

При фиксированном $x$ лагранжиан аффинен по множителям, а инфимум аффинных функций вогнут. Поэтому двойственная задача всегда — максимизация вогнутой функции, как бы ни выглядели $f$ и ограничения.

\begin{theorembox}{Теорема: Слабая двойственность.}
Для любой допустимой $x$, любых $\lambda\ge0$ и $\nu$ выполнено $q(\lambda,\nu)\le f(x)$. Значит, $d^*\le p^*$.
\end{theorembox}

\begin{proofbox}
$q(\lambda,\nu)\le L(x,\lambda,\nu)=f(x)+\sum_i\lambda_ig_i(x)+\sum_j\nu_jh_j(x)\le f(x)$, так как $h(x)=0$ и каждое $\lambda_ig_i(x)\le0$.
\end{proofbox}

\begin{theorembox}{Теорема: Сильная двойственность при условии Слейтера.}
Если задача выпукла, равенства аффинны, $p^*$ конечно и некоторая допустимая точка выполняет все неравенства строго, то $d^*=p^*$ и максимум двойственной задачи достигается.
\end{theorembox}

\begin{whybox}
Считайте $\lambda_i$ штрафом за единицу нарушения $g_i\le0$. Двойственная задача ищет такие штрафы, при которых безусловный минимум оштрафованной цели уже допустим. При слишком низком штрафе минимум жульничает; правильный штраф — это рыночная цена ограничения.
\end{whybox}

\begin{theorembox}{Утверждение: Множители измеряют чувствительность.}
Запишем ограничения как $g_i(x)\le b_i$, и пусть $p^*(b)$ — оптимальное значение. Там, где $p^*$ дифференцируема, а оптимальные множители единственны, $\partial p^*/\partial b_i=-\lambda_i^*$.
\end{theorembox}

\begin{trapbox}
Следите за направлением неравенства. В задаче B1 ограничение — нижняя граница $x+y\ge b$, записанная как $b-x-y\le0$; рост $b$ его ужесточает, и наклон $p^*$ равен $+\lambda^*$, а не $-\lambda^*$.
\end{trapbox}

\problemhead{B1}{Двойственная функция вручную}{У доски}
\textbf{Условие.} Минимизируйте $x^2+3y^2$ при $x+y\ge4$. Запишите ограничение в виде $g\le0$, найдите двойственную функцию $q(\lambda)$, максимизируйте её и сравните $d^*$ с $p^*$. Затем считайте правую часть $b$ параметром и найдите $dp^*/db$ при $b=4$. Укажите $\lambda^*$, $d^*$ и $dp^*/db$.

\textbf{Решение.}
\begingroup\sloppy
\begin{enumerate}[label=\textbf{\arabic*.}, leftmargin=*]
  \item $g=b-x-y\le0$ и $L=x^2+3y^2+\lambda(b-x-y)$. Из $2x-\lambda=0$ и $6y-\lambda=0$ минимум по $(x,y)$ достигается при $x=\lambda/2$, $y=\lambda/6$.
  \item Подставляем: $q(\lambda)=\lambda^2/4+3(\lambda/6)^2+\lambda(4-2\lambda/3)=4\lambda-\lambda^2/3$ — вогнутая парабола. Её максимум при $\lambda\ge0$ там, где $q'(\lambda)=4-2\lambda/3=0$: $\lambda^*=6$, $d^*=12$.
  \item Прямая точка получается из тех же формул: $(x^*,y^*)=(3;\ 1)$; она удовлетворяет $x+y=4$ и даёт $p^*=12=d^*$. Зазора нет, как и обещает условие Слейтера: точка $(4,4)$ строго допустима.
  \item Для произвольного $b$ те же шаги дают $\lambda^*=3b/2$ и $p^*(b)=0{,}75\,b^2$, так что при $b=4$ имеем $dp^*/db=6$ — ровно $\lambda^*$. Знак плюс, потому что $x+y\ge b$ — это $-x-y\le-b$: рост $b$ ужесточает ограничение, и каждая единица стоит $\lambda^*$.
\end{enumerate}
\endgroup
\begin{answerbox}
$\lambda^*=6$, $d^*=p^*=12$, $dp^*/db=6$.
\end{answerbox}

\problemhead{B2}{Сколько стоит ещё единица ресурса}{Самостоятельно}
\textbf{Условие.} Минимизируйте $(x-4)^2+(y-3)^2$ при $x+2y\le b$, $b=3$. Найдите $q(\lambda)$ и точку её максимума, затем решение прямой задачи. Предскажите $p^*$ при $b+0{,}5$ по множителю, а потом вычислите его точно. Укажите $\lambda^*$, $p^*$ и точное $p^*$ при новом $b$.

\textbf{Решение.}
\begingroup\sloppy
\begin{enumerate}[label=\textbf{\arabic*.}, leftmargin=*]
  \item При $c=(4;\ 3)$ и $a=(1;\ 2)$ имеем $L=\|z-c\|^2+\lambda(a^\top z-b)$, где $z=(x,y)$. Приравняв градиент по $z$ к нулю, получим $z(\lambda)=c-\tfrac{\lambda}{2}a$, то есть $x=4-\lambda/2$, $y=3-\lambda$. В самом центре $a^\top c$ сильно превышает $b=3$, так что ограничение окажется точным.
  \item Подстановка $z(\lambda)$ даёт двойственную функцию $q(\lambda)=\lambda(a^\top c-b)-\tfrac{1}{4}\|a\|^2\lambda^2$ — параболу ветвями вниз. Её вершина находится при $\lambda^*=2{,}8$, значение там $9{,}8$. То же $\lambda$, подставленное в $z(\lambda)$, даёт $x^*=(2{,}6;\ 0{,}2)$, для которого $a^\top x^*=b$, и $p^*=\|x^*-c\|^2=9{,}8$: значения совпадают, как и обещает условие Слейтера (начало координат строго допустимо).
  \item Ограничение — это верхняя граница $a^\top z\le b$, поэтому $dp^*/db=-\lambda^*$. При $\Delta b=0{,}5$ прогноз первого порядка равен $p^*-\lambda^*\Delta b=8{,}4$.
  \item Повторное решение с увеличенной правой частью даёт $\lambda=2{,}6$, $x=(2{,}7;\ 0{,}4)$ и $p^*=8{,}45$. Прогноз ошибается на $0{,}05$ — это эффект второго порядка: здесь $p^*(b)=(a^\top c-b)^2/\|a\|^2$ — парабола по $b$, и касательная лежит под ней.
\end{enumerate}
\endgroup
\begin{answerbox}
$\lambda^*=2{,}8$, $p^*=d^*=9{,}8$; после изменения $b$ получаем $p^*=8{,}45$ (прогноз $8{,}4$).
\end{answerbox}

\problemhead{B3}{Равенство и свободный множитель}{Самостоятельно}
\textbf{Условие.} Минимизируйте $x^2+2y^2$ при $h=x+y-6=0$. Найдите двойственную функцию $q(\nu)$, максимизируйте её по всем $\nu\in\mathbb R$ и восстановите прямую точку. Укажите $\nu^*$, $d^*$ и $x^*$.

\textbf{Решение.}
\begingroup\sloppy
\begin{enumerate}[label=\textbf{\arabic*.}, leftmargin=*]
  \item $L=x^2+2y^2+\nu(x+y-6)$. Минимум по $(x,y)$ достигается при $x=-0{,}5\nu$, $y=-0{,}25\nu$.
  \item После подстановки $q(\nu)=-6\nu-0{,}375\nu^2$ — парабола ветвями вниз. Её вершина при $\nu^*=-8$, где $d^*=24$.
  \item Тогда $x^*=(4;\ 2)$ лежит на прямой и $p^*=24=d^*$. Множитель получился отрицательным, и для равенства это допустимо: он лишь показывает, с какой стороны прямой $f$ хотела бы оказаться.
\end{enumerate}
\endgroup
\begin{answerbox}
$\nu^*=-8$, $d^*=p^*=24$, $x^*=(4;\ 2)$.
\end{answerbox}

\begin{selfcheckbox}
Может ли $q(\lambda)$ оказаться выше $p^*$?
\selfcheckanswer Нет, ни при каком $\lambda\ge0$: это слабая двойственность, и ей не нужны ни выпуклость, ни условие Слейтера.
\end{selfcheckbox}

% =====================================================================
\section{От текста задачи к целочисленной модели}\label{sec:modelling}

Чтобы превратить историю в модель, нужны три прохода. Назовите решения и дайте каждому переменную, указав, непрерывная она, целая или двоичная. Переведите каждое правило из текста в ограничение. Наконец запишите цель как функцию переменных. Большинство ошибок — это правило, которое так и не записали.

\begin{definitionbox}{Определение: Целочисленная задача и её релаксация.}
Смешанная целочисленная линейная задача имеет линейную цель, линейные ограничения и часть переменных, обязанных быть целыми. Её линейная релаксация — та же задача без требования целочисленности.
\end{definitionbox}

\begin{theorembox}{Утверждение: Релаксация даёт оценку.}
В задаче максимизации оптимум релаксации не меньше целочисленного оптимума, в задаче минимизации — не больше. Если оптимум релаксации оказался целым, он решает и целочисленную задачу.
\end{theorembox}

\begin{proofbox}
Каждая допустимая целочисленная точка допустима и для релаксации, так что релаксация оптимизирует по более широкому множеству.
\end{proofbox}

\handoutfigure{0.62}{fig_c1}{Релаксация задачи C1 достигает максимума в дробной вершине; лучшая целая точка не получается её округлением.}

\problemhead{C1}{Стулья и столы}{У доски}
\textbf{Условие.} Мастерская делает стулья и столы. Стул приносит прибыль $5$, стол — $4$. На стул уходит $1$ час работы и $4$ единицы древесины, на стол — $3$ часа и $1$ единица. Есть $10$ часов и $13$ единиц древесины. Запишите целочисленную модель, решите её LP-релаксацию, округлите решение LP и найдите целочисленный оптимум. Укажите $z_{LP}$, $z_{IP}$ и число стульев в целочисленном оптимуме.

\textbf{Решение.}
\begingroup\sloppy
\begin{enumerate}[label=\textbf{\arabic*.}, leftmargin=*]
  \item Пусть $x$ — число стульев, $y$ — число столов. Максимизируем $5x+4y$ при $x+3y\le10$ (работа), $4x+y\le13$ (древесина), $x,y\ge0$ и $x,y\in\mathbb Z$.
  \item Отбросим целочисленность. В вершинах $(0,0)$, $(3{,}25,0)$, $\bar x$ и $(0,3{,}3333)$ цель равна $(0;\ 16{,}25;\ 23;\ 13{,}3333)$; лучшая — вершина, где кончаются оба ресурса: $\bar x=(2{,}6364;\ 2{,}4545)$, $z_{LP}=23$.
  \item Округление вниз даёт $(2;\ 2)$ с прибылью $18$. Ближайшая точка решётки $(3;\ 2)$ требует $14$ единиц древесины — больше, чем есть.
  \item Перебор $13$ допустимых точек решётки даёт целочисленный оптимум $(3;\ 1)$ с $z_{IP}=19$: три стула и один стол, а вовсе не округлённый план. Релаксация обещала $23$, так что зазор равен $4$, и ни один целочисленный план эту оценку не превзойдёт.
\end{enumerate}
\endgroup
\begin{answerbox}
$z_{LP}=23$ в точке $\bar x=(2{,}6364;\ 2{,}4545)$; $z_{IP}=19$ в точке $(3;\ 1)$.
\end{answerbox}

\begin{itemize}
  \item \textbf{Включение с постоянными затратами.} $x\le Mz$, $z\in\{0,1\}$, затраты $Fz$: любой выпуск включает эти затраты.
  \item \textbf{Если A, то B.} $z_A\le z_B$ для двоичных индикаторов.
  \item \textbf{Не больше k вариантов из n.} $\sum_iz_i\le k$.
  \item \textbf{Либо одно ограничение, либо другое.} $a^\top x\le b+M(1-z)$ и $c^\top x\le d+Mz$.
  \item \textbf{Произведение двоичных.} $w=z_1z_2$ заменяется на $w\le z_1$, $w\le z_2$, $w\ge z_1+z_2-1$.
\end{itemize}

\begin{trapbox}
Округление оптимума релаксации — не способ решения. Округлённая точка может нарушить ограничение, а даже если она допустима, настоящий целочисленный оптимум может оказаться в нескольких шагах от неё, как в C1.
\end{trapbox}

\begin{notebox}{Каким должно быть M}
Допустимо любое $M$, не меньшее наибольшего значения, которое может принять переменная. Большие значения сохраняют те же целочисленные решения, но ослабляют релаксацию, замедляют метод ветвей и границ, а при огромном $M$ позволяют солверу считать выключенным $z$, отличающееся от нуля на величину допуска.
\end{notebox}

\problemhead{C2}{Включать машину или нет}{Самостоятельно}
\textbf{Условие.} Заказ из $D=7$ единиц можно сделать на двух машинах. Включение машины $i$ стоит $F_i$, единица продукции — $c_i$, и машина делает не больше $u_i$ единиц: $F=(20;\ 6)$, $c=(2;\ 5)$, $u=(5;\ 8)$. Запишите модель с двоичными $z_i$ и $x_i\le M_iz_i$ и найдите самый дешёвый план. Затем замените $z_i\in\{0,1\}$ на $0\le z_i\le1$ и решите релаксацию при $M_i=u_i$ и при $M_i=10\,u_i$. Укажите три значения.

\textbf{Решение.}
\begingroup\sloppy
\begin{enumerate}[label=\textbf{\arabic*.}, leftmargin=*]
  \item Решения: выпуск $x_i\ge0$ и переключатель $z_i\in\{0,1\}$ для каждой машины. Цель $\sum_i(F_iz_i+c_ix_i)$; спрос $x_1+x_2=D$; мощность $x_i\le u_i$; связь $x_i\le M_iz_i$, чтобы ненулевой выпуск включал постоянные затраты. Наименьшее допустимое значение — $M_i=u_i$, и тогда связь уже содержит ограничение мощности.
  \item Вариантов включения всего четыре. Ничего не включать или включить только машину 1 мощностью $5< D$ нельзя: спрос не покрыт. Одна машина 2 стоит $F_2+c_2D=41$. Обе машины, когда дешёвая машина 1 загружена полностью, а остаток делает машина 2, стоят $46$. Целочисленный оптимум — $41$: работает одна машина, причём та, у которой единица дороже.
  \item Ослабим $z_i\in[0,1]$. Теперь каждая единица на машине $i$ стоит $c_i+F_i/M_i$. При $M=u$ эти ставки равны $(6;\ 5{,}75)$, весь спрос уходит на машину 2, и оценка равна $40{,}25$ — чуть ниже целочисленного оптимума.
  \item Возьмём вместо этого $M_i$, равное $10$ мощностям. Ставки станут $(2{,}4;\ 5{,}075)$: постоянные затраты почти исчезли, машина 1 загружена полностью, и оценка падает до $22{,}15$. Целочисленная модель та же самая, но релаксация теперь мало о ней говорит, и методу ветвей и границ пришлось бы работать гораздо дольше.
\end{enumerate}
\endgroup
\begin{answerbox}
Целочисленный оптимум $41$ (работает только машина 2); релаксация $40{,}25$ при $M=u$ и $22{,}15$ при раздутом $M$.
\end{answerbox}

\problemhead{C3}{Рюкзак и его релаксация}{Самостоятельно}
\textbf{Условие.} У четырёх предметов ценности $(8;\ 6;\ 5;\ 4)$ и веса $(3;\ 5;\ 2;\ 8)$; вместимость рюкзака $9$. Запишите модель с двоичными $x_i$. Решите LP-релаксацию, упорядочив предметы по ценности на единицу веса, затем найдите лучшую укладку. Укажите оценку LP, целочисленный оптимум и число уложенных предметов.

\textbf{Решение.}
\begingroup\sloppy
\begin{enumerate}[label=\textbf{\arabic*.}, leftmargin=*]
  \item Модель: $z_i\in\{0,1\}$, максимизировать $\sum_iv_iz_i$ при $\sum_iw_iz_i\le W$, где $v=(8;\ 6;\ 5;\ 4)$, $w=(3;\ 5;\ 2;\ 8)$, $W=9$. По отношению $v_i/w_i$ порядок такой: $(1;\ 3;\ 2;\ 4)$.
  \item Релаксация берёт предметы целиком в этом порядке, пока они помещаются, и долю $0{,}8$ следующего: оценка равна $17{,}8$.
  \item Перебрав помещающиеся наборы, находим лучший — $(1;\ 2)$ ценностью $14$; зазор до оценки равен $3{,}8$. Жадный префикс не обязательно оказывается целочисленным оптимумом.
\end{enumerate}
\endgroup
\begin{answerbox}
Релаксация $17{,}8$, целочисленный оптимум $14$, предметы $(1;\ 2)$.
\end{answerbox}

\begin{selfcheckbox}
Двоичная переменная $z$ входит только в $x\le Mz$ и в затраты $Fz$, $F>0$. Какое значение даёт ей релаксация?
\selfcheckanswer Наименьшее разрешённое, $z=x/M$: релаксация платит только долю $Fx/M$ постоянных затрат.
\end{selfcheckbox}

\end{document}
